\chapter*{符号表}
\phantomsection\addcontentsline{toc}{chapter}{符号表}\label{symbols}
\markboth{符号表}{符号表}
\newcommand{\SymbolPage}[1]{\hyperref[source:\number\numexpr#1+8\relax]{\pageref*{source:\number\numexpr#1+8\relax}}}
\setlength{\LTpre}{0pt}
\begin{longtable}{@{}p{0.32\textwidth}p{0.54\textwidth}r@{}}
\toprule
符号 & 含义 & 页码\\
\midrule
\endfirsthead
\toprule
符号 & 含义 & 页码\\
\midrule
\endhead
\bottomrule
\endfoot
$\sgnmark_n$ & 符号约定中经常出现的项 & \SymbolPage{8}\\[4pt]
$\operatorname{Aut}(M)$ & 正合辛自同构群 & \SymbolPage{96}\\[4pt]
$\operatorname{Aut}(\widehat M)$ & 接上锥体后的辛自同构群 & \SymbolPage{98}\\[4pt]
$\operatorname{Aut}(M,\partial M)$ & 相对于边界的正合辛自同构群 & \SymbolPage{96}\\[4pt]
$\operatorname{Aut}^{\mathrm{gr}}(M,\partial M)$ & 分次辛自同构群 & \SymbolPage{188}\\[4pt]
$\operatorname{Aut}^{\#}(M,\partial M)$ & 膜自同构群 & \SymbolPage{188}\\[4pt]
$CF^{\mathrm{pr}}(L_0,L_1)$ & Floer 上链群，初步版本 & \SymbolPage{109}\\[4pt]
$CF^*(L_0^\#,L_1^\#)$ & Floer 上链群 & \SymbolPage{181}\\[4pt]
$\operatorname{Crit}(\pi)$ & 临界点集 & \SymbolPage{212}\\[4pt]
$\operatorname{Critv}(\pi)$ & 临界值集 & \SymbolPage{212}\\[4pt]
$D(\A),\ D^\pi(\A)$ & 导出范畴 & \SymbolPage{42}\\[4pt]
$D_{S,u}$ & 线性化算子 & \SymbolPage{107}\\[4pt]
$D_{\mathcal S,r,u}$ & 扩展线性化算子 & \SymbolPage{122}\\[4pt]
$\partial^h E$ & 水平边界 & \SymbolPage{212}\\[4pt]
$\partial^v E$ & 垂直边界 & \SymbolPage{212}\\[4pt]
$\operatorname{Fakev}(\Psi,\varpi)$ & 伪临界值集 & \SymbolPage{220}\\[4pt]
$\F(M)^{\mathrm{pr}}$ & Fukaya 范畴，初步版本 & \SymbolPage{126}\\[4pt]
$\F(\pi)$ & Lefschetz 纤维化的 Fukaya 范畴 & \SymbolPage{282}\\[4pt]
$H(\A),\ H^0(\A)$ & 上同调范畴 & \SymbolPage{8}\\[4pt]
$H\F(M)$ & Donaldson--Fukaya 范畴 & \SymbolPage{181}\\[4pt]
$H\F(M)^{\mathrm{pr}}$ & Donaldson--Fukaya 范畴，初步版本 & \SymbolPage{101}\\[4pt]
$HF^{\mathrm{pr}}$ & Floer 上同调，初步版本 & \SymbolPage{101}, \SymbolPage{109}\\[4pt]
$HF^*(L_0^\#,L_1^\#)$ & Floer 上同调 & \SymbolPage{182}\\[4pt]
$\M^{d+1}(y_0,\ldots,y_d)$ & 伪全纯多边形的模空间 & \SymbolPage{126}\\[4pt]
$\M_{E/S}(\{y_\zeta\})$ & 非齐次伪全纯截面的模空间 & \SymbolPage{240}\\[4pt]
$\M_S(\{y_\zeta\}),\ \M_{\mathcal S}(\{y_\zeta\})$ & 伪全纯映射的模空间 & \SymbolPage{105}, \SymbolPage{122}\\[4pt]
$\M_Z^*(y_0,y_1)$ & Floer 轨线的模空间 & \SymbolPage{105}\\[4pt]
$\nufun(\A,\B),\ \operatorname{fun}(\A,\B)$ & $\Ai$ 函子的范畴 & \SymbolPage{10}, \SymbolPage{28}\\[4pt]
$\operatorname{nu\text{-}mod}(\A),\ \modcat(\A)$ & $\Ai$ 模的范畴 & \SymbolPage{19}, \SymbolPage{30}\\[4pt]
$\mathcal R^{d+1},\ \overline{\mathcal R}^{d+1}$ & 带标点圆盘的模空间及其紧化 & \SymbolPage{115}, \SymbolPage{118}\\[4pt]
$TE^v,\ TE^h$ & 沿纤维的切空间及其补空间 & \SymbolPage{212}\\[4pt]
$T_{Y_0}(Y_1)$ & 对象的扭转 & \SymbolPage{63}\\[4pt]
$Z$ & 无限条带 & \SymbolPage{100}\\[4pt]
$Z^\pm$ & 半无限条带 & \SymbolPage{103}\\
\end{longtable}
